% Source: https://artofproblemsolving.com/community/c3h2593709_amc_10_problems
% Stable identity: aops-topic-2593709-q05
Use coordinates $A=(0,0)$, $C=(20,0)$ and, since $\angle ACD=90^\circ$, $D=(20,-30)$. Then $E=(5,0)$, and $B$ lies on line $DE$ on the other side of $AC$ from $D$ (the diagonals of $ABCD$ cross at $E$). Line $DE$ has slope $\dfrac{0-(-30)}{5-20}=-2$, so it is $y=-2(x-5)$.

Since $\angle ABC=90^\circ$, $B$ lies on the circle with diameter $AC$: $(x-10)^2+y^2=100$. Substituting,
\[(x-10)^2+4(x-5)^2=100\iff5x^2-60x+100=0\iff x^2-12x+20=0,\]
so $x=2$ or $x=10$. The point with $x=10$ is $(10,-10)$, on the same side as $D$; the other is $B=(2,6)$.

Hence
\[[ABCD]=[ABC]+[ACD]=\dfrac12\cdot20\cdot6+\dfrac12\cdot20\cdot30=60+300=\boxed{\textbf{(D) }360}.\]

\begin{asy}
usepackage("amsmath");
size(5cm);
defaultpen(fontsize(10pt));
pair A=(0,0),C=(20,0),D=(20,-30),B=(2,6),Ee=(5,0);
draw(A--B--C--D--cycle,linewidth(0.8));draw(A--C);draw(B--D);
draw(circle((10,0),10),gray+dotted);
draw((20,0)--(18.5,0)--(18.5,-1.5)--(20,-1.5));
pair u=unit(A-B),v=unit(C-B);draw(B+1.5*u--B+1.5*u+1.5*v--B+1.5*v);
dot(Ee);
label("$A$",A,W);label("$B$",B,N);label("$C$",C,NE);label("$D$",D,SE);label("$E$",Ee,NE);
label("$20$",(12,0),N);label("$30$",(20,-15),E);label("$5$",(2.5,0),S);
\end{asy}

The dotted circle has diameter $AC$, so it passes through $B$.
